%% =====================================================================
%%  CHAPTER 5: MARKET RESEARCH AND EXISTING PRODUCTS    (budget 6 pages)
%%  chapters/ch5-market.tex   Every number below carries its source;
%%  no unsourced "assumption scores" remain (planning assumptions are labelled).
%% =====================================================================
\pcChapter{market}{6}{חקר שוק ומוצרים קיימים}
\begin{pcguide}[הכסף מדבר]
צריך לשכנע את הקורא שיש כאן “בשר”. לוקחים דוחות שוק שחברות מחקר פרסמו ומחלצים מהם נתונים בצורה מקצועית.
\textbf{לכל מספר, נתון, גרף או נוסחה יש מקור} (\texttt{\textbackslash{}figsrc\{..\}}); ההשוואה היא טכנית ולא שיווקית.
תקציב: \latin{1.5} + \latin{1.5} + \latin{1} + \latin{1} + \latin{1} עמודים לפי הסעיפים.
\end{pcguide}

\section{היקף השוק: \latin{TAM} ו-\latin{SAM}}\label{sec:tam}
שוק הייחוס הרחב הוא תוכנות ניהול ידע ארגוניות, ואילו סימולציות \latin{3D} הן שכבת הדרכה ייעודית מעליהן. שני שווקים אלה
חופפים חלקית, ולכן שוויים אינו מסוכם. לפי \latin{Grand View Research}, שוק תוכנות ניהול הידע העולמי עמד על \money{\$}{20.1B}
בשנת \latin{2024}, והתחזית לשנת \latin{2033} היא \money{\$}{62.2B}, עם שיעור צמיחה שנתי מורכב \latin{(CAGR)} של \pct{13.6}
לתקופה \rng{2025}{2033} \cite{gvr2025market}. בשוק הסמוך של הדרכה אימרסיבית (\latin{VR/AR/MR}) מדווחים אותם אנליסטים על
\money{\$}{16.4B} בשנת \latin{2024} ועל \money{\$}{69.59B} בשנת \latin{2030}, עם \latin{CAGR} של \pct{28.3} לתקופה \rng{2025}{2030}
\cite{gvr2025immersive}. אלה נתוני אנליסט חיצוניים על קטגוריות רחבות: הם מתארים את סביבת הפעולה, ואינם מעידים
על ביקוש מוכח לפתרון הפיילוט.

את \latin{SAM} אפשר לחשב מנתוני החיתוך שפורסמו באותו דוח: ארגונים גדולים מהווים \pct{66.61} מהכנסות שוק תוכנות ניהול הידע,
ופריסת ענן מהווה \pct{65.50} מהן \cite{gvr2025market}, ולכן פריסה שאינה בענן מהווה כ-\pct{34.5}. פיילוט מקומי ומנותק מתאים לארגון
גדול שאינו יכול לשלוח ידע החוצה, ולכן
\begin{equation}
  SAM_{2024}=TAM_{2024}\times s_{\mathrm{large}}\times s_{\mathrm{non\text{-}cloud}}=20.1\times0.6661\times0.345\approx 4.62\ \mathrm{B\$}.
\end{equation}
החישוב מניח שחלקה של הפריסה הלא-ענן בארגונים גדולים דומה לחלקה בכלל השוק; אין בדוח חיתוך משולב, והנחה זו מסומנת כהנחה. בשנת
\latin{2033}, בהנחת שיעורי חלוקה קבועים, ערך זה יגיע לכ-\money{\$}{14.3B}. טבלה~\ref{tab:tam} מרכזת את השכבות, חישובן ומקורן.
\begin{table}[!htbp]
\centering
\caption[שכבות השוק: ערך, חישוב ומקור]{שכבות השוק: ערך, חישוב ומקור (נתוני אנליסט: \latin{Grand View Research} \cite{gvr2025market,gvr2025immersive}; נתוני \latin{IAI}: \cite{iai2025about}; מחיר: \cite{manualto2026pricing})}\label{tab:tam}
\begin{pctable}{c,H{0.25},C{0.16},H{0.25},H{0.22}}
  \pchead{שכבה & הגדרה & ערך & חישוב & מקור ומעמד}
  \pcrow{\latin{TAM} & תוכנות ניהול ידע, העולם, \latin{2024}; תחזית \latin{2033} & \money{\$}{20.1B}; \money{\$}{62.2B} & נתון מדווח; \latin{CAGR} \pct{13.6} (\rng{2025}{2033}) & אנליסט \cite{gvr2025market}}
  \pcrow{\latin{TAM} סמוך & הדרכה אימרסיבית, \latin{2024}; תחזית \latin{2030} & \money{\$}{16.4B}; \money{\$}{69.59B} & נתון מדווח; \latin{CAGR} \pct{28.3} (\rng{2025}{2030}); לא מסוכם עם הקודם & אנליסט \cite{gvr2025immersive}}
  \pcrow{\latin{SAM} & ארגונים גדולים עם פריסה שאינה בענן, \latin{2024} & \money{\$}{4.62B} & \mth{20.1\times0.6661\times0.345} & נתוני חיתוך מדווחים \cite{gvr2025market}; הנחת חלק זהה לארגונים גדולים}
  \pcrow{\latin{SOM} & ראו סעיף~\ref{sec:som} & \money{\$}{76.5K} עד \money{\$}{306K} לשנה & היקף משתמשים \mth{\times} מחיר מקביל & \latin{IAI} \cite{iai2025about}; מחיר \cite{manualto2026pricing}; חדירה: הנחת תכנון}
\end{pctable}
\end{table}

\section{סקירת מוצרים קיימים}\label{sec:products}
נבחרו מוצרים מסחריים המייצגים חמש גישות: הוראות עבודה דיגיטליות, חיפוש ידע ארגוני, לכידת ידע מומחים והדרכה
במציאות רבודה וב\threeD{}, פלטפורמת יישומי שטח, והאלטרנטיבה הידנית. טבלה~\ref{tab:products} מציגה לכל מוצר את מה שניתן
לאמת במקורות פומביים: מודל רישוי ומחיר, מדד ביצועים או זמן הכנה שפורסם, והערה על הפער ביחס לפרויקט. ככל ששדה אינו
מופיע במקור, נכתב ש“לא פורסם”; אין בכך קביעה שהיכולת חסרה, אלא שאין לה ראיה במקורות שנסקרו.
\begin{table}[!htbp]
\centering
\caption[מוצרים מסחריים: רישוי, מדדים שפורסמו ופער]{מוצרים מסחריים: רישוי, מדדי ביצועים שפורסמו והפער ביחס לפרויקט. מקורות: דפי היצרנים \cite{dozuki,poka,glean,guru,vuforia,microsoftguides}; מחירים ומדדים של צד שלישי \cite{manualto2026pricing,manualto2026best,g2dozuki,capterraguru,g2vuforiaec}; נבדקו ב-\latin{26-11-2026}}\label{tab:products}
\begin{pctable}{H{0.17},H{0.20},H{0.27},H{0.18},H{0.18}}
  \pchead{מוצר (קטגוריה) & רישוי ומחיר & מדד ביצועים או זמן הכנה שפורסם & מגבלה ביחס לפרויקט & מעמד המידע}
  \pcrow{\latin{Dozuki} (הוראות עבודה) & הצעת מחיר בלבד \cite{g2dozuki}; אומדן: כ-\money{\$}{850} לחודש, מינימום \latin{50} משתמשים (כ-\money{\$}{17} למשתמש) \cite{manualto2026pricing} & לא פורסם & עיגון \latin{RAG} וניתוק מלא לא צוינו & אומדן צד שלישי}
  \pcrow{\latin{Poka} (הוראות עבודה) & לפי הצעה, מחיר לא פורסם \cite{poka} & לא פורסם & מחולל שאלות מאושר ועיגון בגרסה לא צוינו & דף יצרן}
  \pcrow{\latin{L2L} (הוראות עבודה) & מחיר מפורסם: \money{\$}{150} למשתמש לחודש \cite{manualto2026pricing} & לא פורסם & תרגול \latin{3D}, \latin{RAG} ופריסה מקומית לא צוינו & דיווח צד שלישי}
  \pcrow{\latin{VKS} (הוראות עבודה) & כ-\money{\$}{350} למשתמש לחודש, לפי דיווחי משתמשים \cite{manualto2026pricing}; בקרת גרסאות מובנית \cite{manualto2026best} & לא פורסם & עיגון במקור ותרגול לא צוינו & לא מאומת}
  \pcrow{\latin{Tulip} (יישומי שטח) & תוכנית חינמית; מחיר תוכניות בתשלום לא פורסם \cite{manualto2026pricing} & טענת היצרן: חיסכון של \gesym{80\%} בזמן פיתוח יישום \cite{manualto2026best} & הפקת הדרכה ממסמכי נוהל מגורסאים לא צוינה & טענת יצרן; לא נבדקה}
  \pcrow{\latin{Guru} (חיפוש ארגוני) & הצעת מחיר מותאמת \cite{capterraguru} & לא פורסם & תרגול ומבחן כשירות לא צוינו & דף יצרן \cite{guru}}
  \pcrow{\latin{Glean} (חיפוש ארגוני) & מחיר לא פורסם \cite{glean} & לא פורסם & תרגול וניתוק מלא לא צוינו & דף יצרן}
  \pcrow{\latin{Vuforia Expert Capture} (לכידת מומחה, \latin{AR}) & לא מופיע מחיר \cite{g2vuforiaec} & לא פורסם & עיגון גרסאות במסמך ואישור מומחה לא צוינו & דף יצרן \cite{vuforia}}
  \pcrow{\latin{Dynamics 365 Guides} (\latin{MR}) & פרישה מתוכננת ב-\latin{31-12-2026} \cite{microsoftguides} & לא פורסם & תלות בספק; אינו מתאים לטווח ארוך & דף מחזור חיים}
  \pcrow{אלטרנטיבת האקסל (ידני) & ללא רישוי נוסף; שעות עבודה ידניות & יימדד כבסיס בפיילוט \latin{(D4)} & אין אחזור אוטומטי, יצירה או תרגול & מדידה עצמית}
\end{pctable}
\end{table}
שני מוצרים מוצגים בעיקר כמקרי ניהול סיכון. \latin{Dynamics 365 Guides} מדגים תלות בספק: לפי דף מחזור חיי המוצר, מועד
פרישתו הוא \latin{31-12-2026} \cite{microsoftguides}, ולכן לא ייבחר כבסיס לטווח ארוך. \latin{Vuforia Expert Capture} הוא המתחרה הקרוב ביותר
לרכיב לכידת ידע המומחים: דף היצרן מציג יכולת לכידה ב-\latin{AR/3D} במודל שירות \latin{SaaS}, והבינה המלאכותית מוזכרת בו לבדיקות איכות
בלבד; דף המוצר אינו מציין בקרת גרסאות או אישור מומחה \cite{vuforia}. אלה בדיוק הרכיבים שהפרויקט מוסיף.

\section{מחיר מול שקיפות ביצועים}\label{sec:pricecmp}
המחירים שפורסמו הם מעטים וחלקם משוערים. איור~\ref{fig:pricebar} מציג את כל המחירים למשתמש לחודש שניתן לגזור ממקורות
פומביים לשלושה מוצרים מטבלה~\ref{tab:products}; לשאר המוצרים אין מחיר מפורסם, ולכן אינם מופיעים כעמודה. קנה המידה לוגריתמי,
משום שהפער בין המחירים הוא פי עשרים ויותר.
\begin{figure}[!htbp]
\centering
\pcFig{price_bar}{%
\pcAutoFit{\begin{pctikz}
\begin{axis}[ybar,ymode=log,width=10cm,height=5.8cm,ymin=5,ymax=1000,bar width=1.2cm,enlarge x limits=0.28,
  symbolic x coords={Dozuki,L2L,VKS},xtick=data,ytick={10,100,1000},yticklabels={10,100,1000},
  ylabel={\tn{דולר למשתמש לחודש (סולם לוגריתמי)}},grid=major,
  legend style={at={(0.5,-0.2)},anchor=north,legend columns=3,font=\scriptsize,draw=none},
  nodes near coords,every node near coord/.append style={font=\scriptsize},
  tick label style={font=\small},label style={font=\small}]
  \addplot[fill=pcWarn!80,draw=pcDark,bar shift=0pt] coordinates {(Dozuki,17)};
  \addplot[fill=pcGood!80,draw=pcDark,bar shift=0pt] coordinates {(L2L,150)};
  \addplot[fill=pcBad!70,draw=pcDark,bar shift=0pt] coordinates {(VKS,350)};
  \legend{\tn{אומדן צד שלישי},\tn{פורסם בידי היצרן},\tn{דיווחי משתמשים (לא מאומת)}}
\end{axis}
\end{pctikz}}}
\caption[מחיר רישוי למשתמש לחודש]{מחיר רישוי למשתמש לחודש לפי מוצר. מקורות: \latin{Dozuki} --- אומדן של \latin{Manual.to} (כ-\money{\$}{850} לחודש חלקי \latin{50} משתמשים) \cite{manualto2026pricing}; \latin{L2L} --- מחיר מפורסם כפי שמובא באותו דף \cite{manualto2026pricing}; \latin{VKS} --- דיווחי משתמשים באותו דף (לא מאומת). \latin{Manual.to} היא מתחרה ולכן המחירים מצריכים אימות מול היצרנים}\label{fig:pricebar}
\end{figure}
מחיר אינו מדד איכות, ולכן נבדק גם מה מפורסם על ביצועים. איור~\ref{fig:disclosure} מסמן לכל אחד מתשעת המוצרים אם
במקורות שנסקרו מופיע ערך מספרי עבור חמישה מדדים: מחיר, זמן הכנה, נאמנות או דיוק, ניהול גרסאות ופריסה מקומית.
רק \latin{5} מתוך \latin{45} התאים מלאים (כ-\pct{11}), ואף מוצר אינו מפרסם מדד נאמנות או פריסה מקומית.
\begin{figure}[!htbp]
\centering
\pcFig{disclosure}{%
\pcAutoFit{\begin{pctikz}[x=1cm,y=0.62cm]
  \node[font=\scriptsize\sffamily,text width=1.7cm,align=center,anchor=south] at (0,0.1) {\tn{מחיר}};
  \node[font=\scriptsize\sffamily,text width=1.7cm,align=center,anchor=south] at (-1.9,0.1) {\tn{זמן הכנה}};
  \node[font=\scriptsize\sffamily,text width=1.7cm,align=center,anchor=south] at (-3.8,0.1) {\tn{נאמנות או דיוק}};
  \node[font=\scriptsize\sffamily,text width=1.7cm,align=center,anchor=south] at (-5.7,0.1) {\tn{ניהול גרסאות}};
  \node[font=\scriptsize\sffamily,text width=1.7cm,align=center,anchor=south] at (-7.6,0.1) {\tn{פריסה מקומית}};
  \foreach \nm/\r/\va/\vb/\vc/\vd/\ve in {Dozuki/1/1/0/0/0/0, Poka/2/0/0/0/0/0, L2L/3/1/0/0/0/0, VKS/4/1/0/0/1/0, Tulip/5/0/1/0/0/0,
                                   Guru/6/0/0/0/0/0, Glean/7/0/0/0/0/0, {Vuforia EC}/8/0/0/0/0/0, {Dynamics 365 Guides}/9/0/0/0/0/0}{
    \node[font=\scriptsize\sffamily,anchor=west] at (0.4,-\r) {\nm};
    \foreach \v [count=\k from 0] in {\va,\vb,\vc,\vd,\ve}{
      \ifnum\v=1 \fill[pcDark] ({-1.9*\k},-\r) circle (0.17);
      \else \draw[pcGrey,thick] ({-1.9*\k},-\r) circle (0.17);\fi}
  }
\end{pctikz}}}
\caption[גילוי מדדים מספריים בידי המוצרים]{גילוי מדדים מספריים בטבלה~\ref{tab:products}: עיגול מלא --- ערך מספרי מופיע במקור שנסקר; עיגול ריק --- לא פורסם. מקורות: מחיר --- \cite{manualto2026pricing,g2dozuki}; זמן הכנה (טענת יצרן, \latin{Tulip}) --- \cite{manualto2026best}; ניהול גרסאות (\latin{VKS}) --- \cite{manualto2026best}; דפי היצרנים \cite{poka,glean,guru,vuforia,microsoftguides}. ריק אינו הוכחה להיעדר יכולת}\label{fig:disclosure}
\end{figure}
מכאן שתרשים פיזור של מחיר מול ביצוע מדוד של מוצרים מסחריים אינו אפשרי כרגע: רק שתי טענות זמן הכנה כמותיות פורסמו, שתיהן של
יצרנים ולא של מדידה עצמאית, ואין ביניהן מדד משותף. זהו הנימוק לבנצ'מרק העצמי בסעיף~\ref{sec:bench}: הוא ימדוד זמן הפקה, נאמנות
וניהול גרסאות באותם מסמכים ובאותן שאלות, ויאפשר לראשונה להציב את ביצועי אב-הטיפוס על ציר אחד מול הבסיס הידני.
את חלופות התכן של הפרויקט עצמו (ולא מוצרי שוק) משווה פרק~\ref{ch:alt}.

\section{\latin{SOM}: שוק שאפשר להשיג}\label{sec:som}
ה-\latin{SOM} מוגדר כאן מלמטה למעלה, מאוכלוסיית משתמשים ומחיר מקביל, ולא כאחוז שרירותי מה-\latin{SAM}. נקודת המוצא היא
פריסה ראשונית ביחידות ההדרכה והרכש ובמסגרת טכנית של חטיבה אחת בתעשייה האווירית, ובשלב ב' --- בכל ארבע חטיבותיה, שהן
כ-\latin{15,000} עובדים \cite{iai2025about}; ארגוני תעשיית ביטחון ותעופה דומים הם הרחבה אפשרית שלא כומתה, משום שאין עליהם
מקור שאומת. הערך חושב ממחיר מקביל של \money{\$}{17} למשתמש לחודש, אומדן הצד השלישי הנמוך מבין המחירים שנמצאו
\cite{manualto2026pricing}. שיעור החדירה בתוך החטיבה הוא הנחת תכנון, ולכן הוא מוצג כטווח (טבלה~\ref{tab:som}).
\begin{table}[!htbp]
\centering
\caption[SOM מלמטה למעלה]{\latin{SOM} מלמטה למעלה. נתונים: \latin{15,000} עובדים בארבע חטיבות \cite{iai2025about}; מחיר \money{\$}{17} למשתמש לחודש \cite{manualto2026pricing}; הנחת חלוקה שווה בין חטיבות ושיעור חדירה: הנחות תכנון. ערך שנתי \mth{=} משתמשים \mth{\times 17 \times 12}}\label{tab:som}
\begin{pctable}{H{0.26},C{0.15},C{0.15},C{0.15},C{0.15},C{0.14}}
  \pchead{תרחיש & חדירה & משתמשים (שלב א')& ערך שנתי (שלב א') & משתמשים (שלב ב') & ערך שנתי (שלב ב')}
  \pcrow{זהיר & \pct{5} & \num{188} & \money{\$}{38K} & \num{750} & \money{\$}{153K}}
  \pcrow{\textbf{בסיס} & \pct{10} & \num{375} & \textbf{\money{\$}{76.5K}} & \num{1,500} & \textbf{\money{\$}{306K}}}
  \pcrow{אופטימי & \pct{20} & \num{750} & \money{\$}{153K} & \num{3,000} & \money{\$}{612K}}
\end{pctable}
\end{table}
בתרחיש הבסיס ערך ה-\latin{SOM} הוא \money{\$}{76.5K} בשנה בשלב א' ו-\money{\$}{306K} בשלב ב', כלומר כ-\pct{0.0017} ו-\pct{0.0066} מה-\latin{SAM},
בהתאמה. הנתון הקטן מכוון: הוא מבטא שוק נגיש ממשי לפיילוט ארגוני, ולא פוטנציאל תיאורטי. הוא אינו הכנסה של קבוצת הסטודנטים אלא
הערכת ערך השוק של פריסה כזו לפי מחירי מוצרים דומים, והרגישות העיקרית שלו היא לשיעור החדירה ולמחיר המשתמש. איור~\ref{fig:somfunnel}
מסכם את שלוש השכבות.
\begin{figure}[!htbp]
\centering
\pcFig{som_funnel}{%
\pcAutoFit{\begin{pctikz}[x=1cm,y=1cm]
  \fill[pcPrimary!18,draw=pcPrimary,thick] (0,0) -- (9,0) -- (8,-1.2) -- (1,-1.2) -- cycle;
  \fill[pcPrimary!40,draw=pcPrimary,thick] (1,-1.2) -- (8,-1.2) -- (6.8,-2.4) -- (2.2,-2.4) -- cycle;
  \fill[pcPrimary!78,draw=pcDark,thick] (2.2,-2.4) -- (6.8,-2.4) -- (5.8,-3.6) -- (3.2,-3.6) -- cycle;
  \node[font=\small\sffamily\bfseries] at (4.5,-0.6) {\latin{TAM}: \latin{\$20.1B} (\latin{2024})};
  \node[font=\small\sffamily\bfseries] at (4.5,-1.8) {\latin{SAM}: \latin{\$4.62B}};
  \node[font=\small\sffamily\bfseries,text=white] at (4.5,-3.0) {\latin{SOM}: \latin{\$76.5K}--\latin{\$306K}};
\end{pctikz}}}
\caption[שלוש שכבות השוק]{שלוש שכבות השוק; גודל הצורות אינו בקנה מידה. מקורות: \latin{TAM} --- \cite{gvr2025market}; \latin{SAM} --- חישוב מנתוני החיתוך שבאותו דוח (טבלה~\ref{tab:tam}); \latin{SOM} --- טבלה~\ref{tab:som} לפי \cite{iai2025about,manualto2026pricing}}\label{fig:somfunnel}
\end{figure}

\subsection{רגישות למחיר המשתמש}\label{sec:somsens}
המחיר של \money{\$}{17} הוא הנמוך שבין שלושה מחירי צד שלישי שנמצאו (טבלה~\ref{tab:products}); אותה אוכלוסיית משתמשים בתרחיש הבסיס מניבה ערך שונה בסדר גודל לפי
המחיר שנבחר. טבלה~\ref{tab:somsens} מציגה את ההפרש, כדי שהקורא יראה שההערכה רגישה למחיר לא פחות מאשר לשיעור החדירה.
\begin{table}[!htbp]
\centering
\caption[רגישות ה-SOM למחיר המשתמש]{רגישות ה-\latin{SOM} למחיר משתמש לחודש, בתרחיש הבסיס (\latin{375} משתמשים בשלב א', \latin{1,500} בשלב ב'). מחירים: \cite{manualto2026pricing}; הם דיווחי צד שלישי שלא אומתו מול היצרנים. ערך שנתי \mth{=} משתמשים \mth{\times} מחיר \mth{\times 12}}\label{tab:somsens}
\begin{pctable}{H{0.28},C{0.16},C{0.18},C{0.19},C{0.19}}
  \pchead{מחיר מקביל (לפי מוצר) & מחיר למשתמש לחודש & מכפיל מול הנמוך & ערך שנתי (שלב א') & ערך שנתי (שלב ב')}
  \pcrow{\latin{Dozuki}, אומדן & \money{\$}{17} & \latin{1.0}\mth{\times} & \money{\$}{76.5K} & \money{\$}{306K}}
  \pcrow{\latin{L2L}, מחיר מפורסם & \money{\$}{150} & \latin{8.8}\mth{\times} & \money{\$}{675K} & \money{\$}{2.7M}}
  \pcrow{\latin{VKS}, דיווחי משתמשים & \money{\$}{350} & \latin{20.6}\mth{\times} & \money{\$}{1.575M} & \money{\$}{6.3M}}
\end{pctable}
\end{table}
הטווח הרחב (פי \latin{20.6} בין הקצוות) הוא הסיבה לכך שהדוח אינו מציג את ה-\latin{SOM} כהכנסה צפויה ואינו בונה עליו החלטת
\latin{Go/No-Go}. שני גורמים נוספים מגבילים את ההערכה: ראשית, ארגון כמו \latin{IAI} עשוי לרכוש רישיון אתר או לפתח פנימית ולא לשלם
לפי משתמש, ולכן מחיר למשתמש הוא קירוב ולא תחזית מסחרית; שנית, שיעור חדירה של \pct{10} הוא הנחת תכנון שלא נבדקה מול
נתוני שימוש של הארגון. השימוש הנכון בטבלה הוא כטווח התייחסות לעלות החלופה של הארגון (מה היה משלם על מוצר מקביל), ולא כמדד להצלחת
הפרויקט, שנמדדת ביעדי פרק~\ref{ch:problem}.

\section{בנצ'מרק}\label{sec:bench}
\begin{pcguide}[בנצ'מרק לעומת חקר שוק]
חקר השוק שואל “מה קיים ובכמה?”; הבנצ'מרק שואל “מי עושה את זה הכי טוב, איך, ובכמה אנחנו רחוקים ממנו?”.
הפער שנמדד הופך ליעד מדיד בדוח, ובהמשך למה שהסימולציה צריכה להראות.
\end{pcguide}
לאחר שהוכח שהשוק אינו מפרסם נתוני ביצוע (איור~\ref{fig:disclosure}), הבנצ'מרק מגדיר מדידה עצמית על אותו סט מסמכים ושאלות לכל
חלופה, בחמישה שלבים (טבלה~\ref{tab:benchmark}). המדדים: חציון דקות ליחידה מאושרת, שיעור טענות נתמכות במדגם בידי שני
מעריכים, מספר פריטים מיושנים שלא נחסמו לאחר שינוי גרסה, מספר בקשות לרשת חיצונית (יעד אפס) ושיעור שאלות שעברו את מחוון האיכות.
לתרחיש ה-\latin{3D} יימדדו בנוסף זמן בניית תרחיש מאושר, אחוז הפעולות שתועדו, שיעור ההשלמה ללא עזרה, וזמן ועלות עדכון הסצנה לאחר
שינוי נוהל. הפער שיימדד יוזן חזרה לספי הדרישות בפרק~\ref{ch:problem}.
\begin{table}[!htbp]
\centering
\caption{חמשת שלבי הבנצ'מרק}\label{tab:benchmark}
\begin{pctable}{c,H{0.22},H{0.40},H{0.24}}
  \pchead{שלב & משימה & מה עושים בפרויקט שלנו & תוצר}
  \pcrow{\num{1} & מה משווים & בחירת סט של \latin{20} שאלות ומסמכים, זהה לכל חלופה & תיק מבחן}
  \pcrow{\num{2} & מול מי & הכנת יחידת לימוד ידנית (בסיס) ועם אב-הטיפוס; מסגרת מאמר הדגל כייחוס; בדיקת דמו של מוצר שיתאפשר לבדוק ללא עלות & שני תוצרים שקולים}
  \pcrow{\num{3} & איסוף נתונים & מדידת זמן, טענות נתמכות ובקרת גרסאות; לכל נתון מקור & קובץ תוצאות}
  \pcrow{\num{4} & מדידת הפער & חישוב הפער מול ספי הדרישות (\latin{REQ}) ומול מוצר הייחוס & דוח פערים}
  \pcrow{\num{5} & יעדים ופעולה & תיקון המודל וחזרה על הניסוי & החלטת \latin{Go/No-Go}}
\end{pctable}
\end{table}
\section{מה השוק מלמד ומה נגזר לתכן}\label{sec:marketimpl}
הממצאים של הפרק מתורגמים להחלטות תכן ומדידה בטבלה~\ref{tab:marketimpl}. כל שורה נשענת על טבלה או איור קודמים בפרק, ולכן אינה מוסיפה טענת שוק חדשה.
\begin{table}[!htbp]
\centering
\pcCompactTables
\caption[השלכות חקר השוק על התכן]{השלכות חקר השוק על התכן. מקורות: טבלאות~\ref{tab:tam}, \ref{tab:products} ו-\ref{tab:somsens} ואיור~\ref{fig:disclosure}, כלומר דפי היצרנים ודיווחי צד שלישי שצוינו שם}\label{tab:marketimpl}
\begin{pctable}{H{0.26},H{0.38},H{0.36}}
  \pchead{ממצא בפרק & משמעות & נגזרת לתכן ולמדידה}
  \pcrow{אין מוצר שנסקר שמפרסם מדדי נאמנות או זמן הפקה מאומתים (איור~\ref{fig:disclosure}) & השוואה כמותית מול השוק אינה אפשרית מנתונים שפורסמו & הבנצ'מרק מודד בעצמו את הבסיס הידני ואת אב-הטיפוס על אותו סט שאלות (סעיף~\ref{sec:bench})}
  \pcrow{טווח המחירים רחב מאוד, והוא דיווח צד שלישי שלא אומת (טבלה~\ref{tab:somsens}) & ערך ה-\latin{SOM} אינו בסיס להחלטה כלכלית & החלטת \latin{Go/No-Go} נשענת על מדדי פרק~\ref{ch:problem} ולא על הערכת הכנסה}
  \pcrow{מוצרים מסוימים מציגים ניהול גרסאות או זמן הכנה מופחת כטענת יצרן (טבלה~\ref{tab:products}) & גרסאות וקיצור זמן הם ציפייה בסיסית ולא חידוש בפני עצמם & הבידול נמדד בשילוב: פריסה מקומית, עיגון לגרסה ותצורת תרחיש \threeD{} שמתעדכנת}
  \pcrow{ה-\latin{SAM} נגזר מנתוני חיתוך כלליים של שוק ניהול ידע (טבלה~\ref{tab:tam}) & אין נתון על שוק ההדרכה ההנדסית בפריסה מקומית & ההיקף המסחרי אינו נטען; הפרויקט נמדד כפיילוט ארגוני}
\end{pctable}
\end{table}
\pcChapterEnd{market}
